\section*{Appendix}
\label{appx:intro}
In this appendix, we provide additional details about the implementation and usage of the \benchmark{} framework and the \texttt{\benchmark{}Env} interface. We also provide visualizations and analyses of additional experiments to demonstrate \benchmark{}'s utility and garner further insight into the extent of current models' performance on the interactive coding task. The full template for each prompting strategy is also included. Finally, we also discuss some of the impacts, risks, and limitations of our work. The webpage for \benchmark{} is \url{https://intercode-benchmark.github.io/}. The code for \benchmark{} is \url{https://github.com/princeton-nlp/intercode}; the link is also included on the \benchmark{} webpage.

\section{Environment Details}
\label{appx:env}

\subsection{\benchmark{} Interface}
\label{appx:env:interface}
The \benchmark{} interface inherits the OpenAI gym~\cite{openai_gym} environment API definition. 
Specifically, \texttt{\benchmark{}Env} is written as an abstract class that primarily handles the main execution logic for processing code interactions, in addition to logging, data management, and sand-boxed execution, along with both environment-level and task-level customization. 

\texttt{\benchmark{}Env} exposes the following API. Creating an interactive coding environment requires defining a subclass of \texttt{\benchmark{}Env}. The methods denoted with an asterisk can be overridden for the purposes of customization.

\texttt{\_\_init\_\_(self, data\_path: str, image\_name: str, **kwargs)}
\begin{itemize}[noitemsep,topsep=0pt]
    \item Validates that the dataset specified by \texttt{data\_path} is formatted correctly and can be used in an interactive setting.
    \item Uses the Docker image specified by \texttt{image\_name} to create and connect with a Docker container instance of the image.
    \item Initializes Logging Handler
    \item Keyword arguments:
    \begin{itemize}[noitemsep]
        \item \texttt{verbose (bool)}: If true, logging is enabled and environment interactions are shown to standard output
        \item \texttt{traj\_dir (str)}: If a valid path is provided, task episode summaries are saved to the given directory (generated by \texttt{save\_trajectory})
        \item \texttt{preprocess (callable)}: If provided, this function is run before every task episode. It is a way to provide task instance-specific customization of the execution environment.
    \end{itemize}
\end{itemize}

\texttt{reset(self, index: int = None) -> Tuple[str, Dict]}
\begin{itemize}[noitemsep,topsep=0pt]
    \item Retrieves task record from data loader
    \item Calls \texttt{reset\_container}
    \item Reset task level logger, instance variables
\end{itemize}

\texttt{step(self, action: str) -> Tuple[str, int, bool, Dict]}
\begin{itemize}[noitemsep,topsep=0pt]
    \item Log (action, observation)
    \item Invoke \texttt{exec\_action} on action argument
    \item If \texttt{action}=\texttt{submit}, invoke \texttt{get\_reward}, \texttt{save\_trajectory}
\end{itemize}

\texttt{save\_trajectory(self)}
\begin{itemize}[noitemsep,topsep=0pt]
    \item Saves task metadata, (action, obs.) sequence, and reward info to \texttt{.json} in \texttt{traj\_dir}
\end{itemize}

\texttt{close(self)}
\begin{itemize}[noitemsep,topsep=0pt]
    \item Safely exit or stop any resources (i.e. docker container) used by the environment
\end{itemize}

\texttt{* execute\_action(self, action: str)}
\begin{itemize}[noitemsep,topsep=0pt]
    \item Defines how the action is executed within the context of the docker container.
    \item Requires impl. because the Dockerfile definition, particularly its \textit{entrypoint}, affects how an action would be invoked within the container.
    \item Default impl. passes the \texttt{action} string directly into a \texttt{self.container.exec(action)} call, which invokes the action in the environment and returns execution output. A timeout is imposed on execution duration.
\end{itemize}

\texttt{* get\_reward(self) -> Tuple[float, Dict]}
\begin{itemize}[noitemsep,topsep=0pt]
    \item Handles reward calculation of actions with respect to the gold command(s) for a task episode.
    \item Requires impl. because the concept and scoring for task completion varies across datasets and environments.
\end{itemize}

\texttt{* reset\_container(self)}
\begin{itemize}[noitemsep,topsep=0pt]
    \item Handles resetting of execution container (i.e. resetting file system to original state).
    \item Requires impl. because the approach to restoring a setting to its initial state varies.
\end{itemize}

\begin{figure}
    \centering
    \includegraphics[width=0.85\textwidth]{figures/assets/interface-visual.pdf}
    \caption{Visualization demonstrating the intended invocations and usage of the \texttt{\benchmark{}Env} interface, along with how the functions requiring implementation (\texttt{get\_reward()}, \texttt{execute\_action()}, \texttt{reset\_container()} are called by the methods of the main interactive loop.}
    \label{appx:fig:interface_visual}
\end{figure}

Figure~\ref{appx:fig:interface_visual} conveys how each of these methods are invoked and how they related to one another. In summary, the technicalities for setting up an interactive coding task for a specific system with one or more programming languages as the action space involve:

\begin{itemize}[noitemsep,topsep=0pt]
    \item Defining a \texttt{Dockerfile}
    \item Providing a dataset with the \texttt{query} and \texttt{gold} fields
    \item (Optional) Defining a reward (\texttt{get\_reward}) function to define task completion.
    \item (Optional) Creating an \texttt{\benchmark{}Env} subclass that overrides the \texttt{execute\_action} and \texttt{get\_reward} methods
\end{itemize}

\subsection{Bash Environment}
\label{appx:env:bashenv}
\textbf{Environment definition.} The \texttt{Dockerfile} defining the Bash-based environment is founded on the LTS version of the Ubuntu operating system. Several Linux dependencies that can potentially be used by an agent to address instructions in the \benchmark{}-Bash Dataset are then installed via the Advanced Package Tool (\texttt{apt}) interface. Next, a shell script is invoked within the \texttt{Dockerfile} to initialize one of the three file systems displayed in Figure~\ref{fig:nl2b_filesystems}. The shell script consists of a simple sequence of \texttt{mkdir}, \texttt{touch}, and \texttt{echo} commands to deterministically create and populate the content of multiple files and folders. Finally, \texttt{git} is configured for the purposes of determining file diffs per task episode (\texttt{git status -s}) and resetting an environment to its original state (\texttt{git reset --hard; git clean -fd;}) before the beginning of a new task episode. The original code for the \texttt{Dockerfile} along with the file system creation scripts can be found on the project GitHub repository.

% \textbf{\benchmark{}-Bash Dataset.} The \benchmark-Bash dataset is built over NL2Bash \cite{lin-etal-2018-nl2bash} which lacks specificity in queries and grounding to any underlying file system, hindering it from being directly used for interactive evaluations. 
% To overcome this limitation, we process NL2Bash as follows:
% \begin{itemize}[noitemsep,topsep=0pt]
% \item Consider a subset of 1000 commands with each having $\geq$ 4 utilities
% \item Filter out commands with unsupported utilities (eg. "ssh", "sudo", time and GUI dependent utilities)
% \item Filter out non-UNIX, non-Linux commands
% \item Enhance under-specified commands with specific file names/directory names/paths
% \item Update deprecated utilities/flags
% \end{itemize}
% The resulting 200 commands were grouped into 4 disjoint sets, 3 of which were grounded to custom-designed file systems, while one set is file-system agnostic. This categorization allowed for a comprehensive evaluation of different command-grounding scenarios.


\textbf{Dataset details.} The log-frequency distribution of the top-50 utilities is displayed in Figure~\ref{fig:nl2b_utility_distribution}.
The NL2Bash~\cite{lin-etal-2018-nl2bash} dataset is made available for use under the \hyperlink{https://www.gnu.org/licenses/gpl-3.0}{GPLv3 License}.
To assess the generalizability of our approach, we designed three distinct file systems to accommodate the bash commands we collected. A key consideration during the construction of these file systems was to ensure that a significant portion of the executed commands would not result in operations that yield no changes. This deliberate design choice aimed to provide a more comprehensive assessment of our approach's adaptability and effectiveness across various scenarios and command executions. The file systems encompass a wide range of file types, including text files (.txt), program files (.c, .java, .py), compressed files (.gz), shell scripts (.sh), PHP scripts (.php), JSON files (.json), documents (.doc), spreadsheets (.csv), webpages (.html), database schemas (.sql), hidden files, and files with special characters in their names, convoluted folder hierarchies. Their directory structures are illustrated in Figure~\ref{fig:nl2b_filesystems}. For simplicity, we consider the top-level folder created within the root directory (\texttt{testbed}, \texttt{system}, \texttt{workspace}) as the root of each file system. This root folder contains files and sub-folders that necessitate access and manipulation, while changes are monitored throughout the entire container to accurately evaluate the models' actions. Notably, we intentionally designed file system 1 to be more intricate and encompass relatively challenging bash tasks compared to the other two file systems. Thereby, the models' performance is relatively lower for file system 1.

\begin{figure}
\begin{subfigure}[t]{0.25\textwidth}
\dirtree{%
.1 /testbed.
.2 FooBar.html.
.2 Hello.java.
.2 Hello1.java.
.2 NewClass.java.
.2 dir1.
.3 AnotherHello.java.
.3 info.php.
.3 subdir1.
.4 jsonfile1.json.
.4 pythonscript4.py.
.4 shellscript1.sh.
.4 subsubdir1.
.5 pythonscript1.py.
.5 shellscript4.sh.
.4 textfile4.txt.
.3 subdir2.
.3 textfile1.txt.
.2 dir2.
.3 shellscript2.sh.
.3 subdir1.
.4 javafile1.java.
.4 textfile2.txt.
.3 subdir2.
.4 pythonscript2.py.
.4 shellscript5.sh.
.4 subsubdir1.
.4 textfile5.txt.
.2 dir3.
.3 subdir1.
.4 pythonscript3.py.
.4 subsubdir1.
.5 FooBar.
.6 file.txt.
.5 file.txt.
.5 shellscript3.sh.
.5 textfile3.txt.
.5 tmp.
.6 tmp.txt.
.3 subdir2.
.4 csvfile1.csv.
.3 textfile6.txt.
.2 hello.c.
.2 hello.php.
.2 recent.txt.
.2 textfile7.txt.
}
\caption{File System 1}
\label{fig:nl2b_fs_1}
\end{subfigure}
\hfill
\hfill
\begin{subfigure}[t]{0.25\textwidth}
\dirtree{%
.1 /system.
.2 .DS\_Store.
.2 MANIFEST.
.2 a.out.
.2 folder1.
.3 a.out.
.3 data.csv.
.3 doc1.doc.
.3 doc2.doc.
.3 keep.txt.
.3 log1.log.
.3 new.sh.
.3 old2.txt.
.3 recent.txt.
.3 script1.sh.
.3 text2.txt.
.3 text3.txt.
.3 text4.txt.
.2 folder2.
.3 special text3.txt.
.3 special\_text1.txt.
.3 special\_text2.txt.
.3 text1.txt.
.2 folder2.tar.gz.
.2 folder3.
.3 backup\_dbg.
.4 backup.
.4 sql1.sql.
.4 text1\_dbg.txt.
.3 special text4.txt.
.3 temp.
.4 empty.txt.
.4 temp1.
.5 temp.
.5 temp\_1.
.5 text1.txt.
.2 html1.html.
.2 temp.
.3 keep.txt.
.2 text1.txt.
.2 text3.txt.
}
\caption{File System 2}
\label{fig:nl2b_fs_2}
\end{subfigure}
\hfill
\begin{subfigure}[t]{0.25\textwidth}
\dirtree{%
.1 /workspace.
.2 .hidden.txt.
.2 archive.tar.gz.
.2 dir1.
.3 .hidden1.txt.
.3 a.txt.
.3 all.txt.
.3 file.c.
.3 file.txt.
.3 hello.c.
.3 hello.txt.
.3 long.txt.
.3 new1.sh.
.3 readonly.txt.
.3 script1.sh.
.3 sum.c.
.3 terminate.txt.
.2 dir2.
.3 csvfile1.csv.
.3 foo.txt.
.3 hello.txt.
.3 mysql.
.4 sql1.sql.
.2 files.txt.
.2 new.sh.
.2 old2.txt.
.2 recent1.txt.
.2 results.txt.
.2 test.
.3 2dir.
.3 3dir.
.3 dir1.
}
\caption{File System 3}
\end{subfigure}
\caption{File System structures designed for \benchmark{}-Bash.}
\label{fig:nl2b_filesystems}
\end{figure}

\textbf{Reward function.} Evaluation of an agent's trajectory across a single task episode towards carrying out the given instruction is determined by modifications to the file system and the latest execution output. The instructions found in the \benchmark{}-Bash dataset fall under one of two buckets: it either 1. Requests information about the file system that can be answered via execution output generated from a correct sequence of Bash actions (i.e. "How many files...", "What is the size of...", "Where is the .png image stored?") or 2. Requests a change to the location, configuration, or content of a file or folder (i.e. "Move the \texttt{dir1} folder from...", "Set the permissions to...", "Append a line to..."). Any relevant correct changes are therefore captured by considering both execution output and file system modifications during evaluation.

We define $A$ and $G$ as the outputs of the \texttt{agent} and \texttt{gold} commands respectively, where $A_{out}$ and $G_{out}$ refer to the execution output, and $A_{fs}$ and $G_{fs}$ refer to a list of entries reflecting file system modifications, where each entry is \texttt{[file path, modification type $\in$ [added, changed, deleted]]}. We then formally define the reward function as follows:

\begin{equation}
    \label{eq:reward_bash}
    \begin{gathered}
        \mathcal{R} = 0.34 * \texttt{similarity}(A_{out}, G_{out}) \\
        + 0.33 * (1 - \texttt{erf}(\lvert A_{fs} \cup G_{fs} - A_{fs} \cap G_{fs} \rvert)) + \\
        + 0.33 * \frac{\texttt{is\_correct}(A_{fs} \cap G_{fs})}{A_{fs} \cap G_{fs}}
    \end{gathered}
\end{equation}

Where \texttt{similarity} refers to lexical similarity, which is determined by the cosine similarity score between TF-IDF vectors (calculated with \texttt{TfidfVectorizer} from \texttt{scikit-learn}) of the two execution outputs. The second component of the reward function reflects the number of file system modifications that were either not completed or not necessary; the error associated with the total number of misses is constrained to the range \texttt{[0,1]} using the Gauss error function (\texttt{erf}), where $0$ corresponds to no file system modification mistakes. The third component checks what proportion of paths altered by both \texttt{agent} and \texttt{gold} were modified correctly. The \texttt{is\_correct} function returns the number of file paths that were changed correctly, determined by checking whether the \texttt{md5sum} hashes of each file path are identical for \texttt{agent} and \texttt{gold}. If $A_{fs} \cap G_{fs} = \emptyset$, this reward is automatically $1$. The scalar weights for each component are arbitrarily assigned.

A max score of $1$ is achieved only if the correct file paths are changed, the changes are correct, and the latest execution output matches the \texttt{gold} command output exactly. Figure~\ref{fig:preview} visualizes the reward function.
While an exact match comparison would have been a simpler choice to satisfy the Success Rate metric put forth in the main paper, we design this reward function to 1. Demonstrate that \benchmark{} can support complex reward functions that account for multiple forms of execution output, and 2. Provide practitioners who use the \benchmark{}-Bash environment with a scalar reward that reflects how "similar" the given output is to the expected, rather than a flat 0/1 reward value that may over-penalize and discount the efforts of more capable reasoning abilities. These reasons also motivate the SQL-based environment's reward function, discussed in the following section.

\begin{figure}
    \centering
    \begin{minipage}{0.49\textwidth}
        \centering
        \includegraphics[width=\textwidth]{figures/assets/bash_util_freq.pdf}
        \caption{Top 30 most frequently occurring bash utilities out of the 66 in \benchmark{}-Bash with their frequencies in log scale. }
        \label{fig:nl2b_utility_distribution}
    \end{minipage}\hfill
    \begin{minipage}{0.49\textwidth}
        \centering
        \includegraphics[width=\textwidth]{figures/assets/sql_hardness.png}
        \caption{Distribution of \texttt{gold} command difficult for \benchmark{}-SQL task data adapted from the Spider SQL dataset.}
        \label{fig:sql_hardness_distribution}
    \end{minipage}
\end{figure}

\subsection{SQL Environment}
\label{appx:env:sqlenv}
\textbf{Environment Definition.} The \texttt{Dockerfile} defining the SQL-based environment inherits from the MySQL image and adds a \texttt{.sql} file setup script to the \texttt{/docker-entrypoint-initdb.d} directory within the Docker image; this is a special directory made for container initialization. On container start-up, the added \texttt{.sql} file, which creates and populates databases with tables and tables with records, is automatically invoked. Since the \benchmark{}-SQL dataset does not feature any queries that involve modifying the database in any manner (i.e. no \texttt{INSERT}, \texttt{UPDATE}, or \texttt{DELETE} commands), there is no reset mechanism written into the \texttt{Dockerfile} definition that is invoked before each task episode; with that said, adding a reset script or version control to the \texttt{Dockerfile} is simple.

\textbf{\benchmark-SQL dataset.}
\benchmark-SQL is adopted from the development set of the Spider dataset \cite{yu-etal-2018-spider}. Spider 1.0 is a large-scale cross-domain dataset on generating SQL queries from natural language questions whose development set contains 1034 pairs of <\texttt{instruction}, \texttt{gold}> task instances spanning 20 databases.
The distribution of queries according to their hardness criterion is shown in Figure~\ref{fig:sql_hardness_distribution}.
As discussed in Section~\ref{sec:benchmark:ic_envs}, a filtering criterion narrows down the Spider dataset's information to only the necessary components.
We do not add anything to the Spider dataset that was not originally available.
The Spider 1.0 dataset is available for use under the \hyperlink{https://creativecommons.org/licenses/by-sa/4.0/legalcode}{CC BY-SA 4.0 license}.

\begin{figure}
    \centering
    \begin{minipage}{0.48\textwidth}
        \centering
        \includegraphics[width=\textwidth]{figures/assets/Traj-SQL}
        \caption{Example of interactions between an agent and the \benchmark{} SQL Environment}
        \label{fig:trajectory_sql}
    \end{minipage}\hfill
    \begin{minipage}{0.48\textwidth}
        \centering
        \includegraphics[width=\textwidth]{figures/assets/Reward-SQL}
        \caption{Evaluation of the results of agent interactions with the SQL Environment against the gold command associated with the task. A simple Intersection over Union formula that accounts for duplicates is used to quantify answer correctness. Task completion is a reward of 1.}
        \label{fig:reward_sql}
    \end{minipage}
\end{figure}

\textbf{MySQL databases.} We first resolve data types for primary, foreign key pairs across the provided table schemas in Spider for conflicting instances and generate the corresponding SQLite databases. Next, to align with our Docker-supported environment, we convert the SQLite databases to MySQL format using \texttt{sqlite3mysql}~\cite{sqlite3mysql}, a Python library, and then generate a unified MySQL dump having schemas for all the tables. To handle case-sensitive table name discrepancies between the queries and the underlying schema in the original Spider dataset, we activate the \texttt{lower\_case\_table\_names} setting in our evaluation environment. Additionally, for proper access controls, we create a test user and grant them all privileges for all the tables.

\textbf{Reward function.} The completion evaluation mechanism compares the output of the \texttt{gold} SQL command with the latest execution output (i.e. latest observation) from the agent's interaction trajectory. The execution output of all \texttt{gold} SQL queries is a list of records. Each record is a tuple of one or more values that may be different types. For any single execution output, the order of types for every record is identical. Given the \texttt{agent} command(s)' latest execution output $A$ and the \texttt{gold} command's execution output $G$, we formulate the reward function as follows:

\begin{equation}
    \label{eq:reward_sql}
    \mathcal{R} = \frac{A \cap G}{A \cup G} * (kendalltau((A \cap (A \cap G)), (G \cap (A \cap G))) + 1) / 2
\end{equation}

We employ Intersection over Union (\textit{IoU}), or more formally the Jaccard Index, to quantify the correctness of the latest execution output generated by the agent against the \texttt{gold} output. If the latest execution output of the SQL query is not in the form of a list of records (i.e. a string error message), the reward is $0$ by default. Among the items that lie in the intersection of the \texttt{agent} and \texttt{gold} execution outputs, we also apply a penalty if the records are in the incorrect order. Since achieving the correct order of fields in a record is of non-trivial importance to addressing many SQL queries correctly, we do not do any re-ordering or pre-processing of the list of records. Therefore, a record formatted as \texttt{("Ross", 29)} is not awarded any credit against a \texttt{gold} output that includes \texttt{(29, "Ross")}. To quantify how sorted the \texttt{agent} output is relative to the \texttt{gold} output, we lean on Kendall's $\tau$ and adjust the output range to $[0, 1]$. The \textit{IoU} score is then directly scaled by this coefficient.

All in all, only a correctly ordered list with the exact set of records found in the \texttt{gold} output would receive a max score of 1, which corresponds to task completion.
Figure~\ref{fig:reward_sql} visualizes the reward function for an example set of outputs. Note that in the main paper, the Success Rate metric is used; the scalar $3/7$ output shown in the figure is treated as a 0 when quantifying whether the task was completed via the 0/1 Success Rate metric. As mentioned in the discussion of the Bash reward function, this reward function also aims to be a richer and fairer continuous evaluation metric of a model's reasoning abilities compared to a binary 0/1 task completion score.

\subsection{Python Environment}
\textbf{Environment definition.} The \benchmark{}-Python task environment inherits from a bare-minimum Python 3.9 image that provides the basic essentials for initializing a Python interpreter.
We were unable to determine how to initialize a Python interpreter within a Dockerfile such that the container would then be capable of automatically executing Python commands sent to it while continuous logging every action/observation per turn.
To overcome this, we create and define a backend application that runs within the Docker container, simulates a Python interpreter, and is responsible for handling input/output.
By having the application sit between the agent's actions and the interpreter, we are able to log every episode faithfully in addition to providing an environment that is agent-friendly and faithful to the experience of a real Python interpreter.

\textbf{\benchmark{}-Python dataset.} A large majority of code datasets popular within the NLP community are based on Python and present code completion as the primary task~\cite{chen2021evaluating,austin2021program,hendrycksapps2021}.
In the original problem setting, a task worker is asked to synthesize code in a zero, one, or few-shot setting with little to no access to an execution environment.
In the interactive setting, task workers are asked to accomplish the same objective, but informed that they have a Python environment to do whatever may help them write the function correctly, such as prototype different implementations and write/execute their own unit tests.
Therefore, datasets such as HumanEval, APPS, and MBPP require little to no revisions to be usable within the \benchmark{} environment, with the only necessary processing for all three being renaming of dataset attributes to \benchmark{}-compatible names.
A visualization of an example trajectory of interactions between an agent and the Python interpreter is presented in Figure~\ref{fig:trajectory_python}.

\textbf{Reward function.} We preserve the original metric of proportion of unit tests passed to evaluate agent implementations, with all tests passing being equivalent to task completion.
Complementary to the visualization of interactions, we also show how \benchmark{}-Python performs automatic evaluation of an agent's implementation of the desired function in Figure~\ref{fig:reward_python}.

\begin{figure}
    \centering
    \begin{minipage}{0.48\textwidth}
        \centering
        \includegraphics[width=\textwidth]{figures/assets/Traj-Py.pdf}
        \caption{Example of interactions between an agent and the \benchmark{} Python Environment}
        \label{fig:trajectory_python}
    \end{minipage}\hfill
    \begin{minipage}{0.48\textwidth}
        \centering
        \includegraphics[width=\textwidth]{figures/assets/Reward-Py.pdf}
        \caption{In this setting, an agent interacts with a Python Interpreter to 1. implement the requested method and 2. write test cases to determine function correctness. Upon submission, the reward function then evaluates the agent's implementation with a set of unit tests.}
        \label{fig:reward_python}
    \end{minipage}
\end{figure}